\section{Generátor cílového kódu (gencode)}

Modul \texttt{codegen} realizuje sestavení cílového kódu ve formátu IFJcode25 z instrukcí tříadresového kódu. Na rozdíl od implementace, kde je cílový kód generován přímo z AST, jsme se rozhodli generovat také vnitřní kód. Mezikód pomohl oddělit generátor od syntaktické a sémantické analýzy, takže bylo pohodlnější soustředit se na generování kontrol běžných chyb (chyb typů) a testování. 

Generátor je implementován ve dvou souborech: 1. \texttt{codegen.h} obsahuje deklaraci veřejných funkcí (\texttt{generate\_code}, \texttt{string\_to\_ifjcode}), 2. \texttt{codegen.c} obsahuje implementace veřejných a lokálních funkcí.

\subsection*{Generování kódu}

Funkce \texttt{generate\_code(...)} je vstupem do generátoru cílového kódu. Prochází seznamem instrukcí a v závislosti na operaci volá potřebnou funkci, přičemž každá operace z TAC má svou vlastní funkci. Každá funkce generuje potřebné kontroly typů a převody int na float nebo float na int, pokud float v desetinné části obsahuje nulu. Tento přístup pomáhá při testování a zvyšuje čitelnost kódu.

Každá funkce operace používá menší funkce pro generování jednotlivých instrukcí, například  \texttt{gen\_defvar}, \texttt{gen\_move}. Může také volat funkce, které generují hotovou sérii instrukcí (\texttt{gen\_type}, \texttt{gen\_same\_operand\_check} atd.) nebo funkce pro generování části instrukce (\texttt{gen\_operand}).

Generátor také vytváří své jedinečné návěští pomocí externí funkce \texttt{create\_unique\_label}, která pokračuje v číslování po generátoru vnitřního kódu.

\subsection*{Použité datové struktury}

Modul používá dvě hlavní datové větve. První větev tvoří instrukce. Jsou implementovány pomocí TacInstruction, který obsahuje datové typy TacOperationCode a Operand. Tyto dvě struktury jsou pro generování klíčové, protože obsahují názvy návěští, identifikátory, konstanty, operace atd. Každá instrukce obsahuje odkaz na seznam TACDLList.

Druhou větví je tabulka symbolů. Pomáhá na začátku generování generovat globální proměnné, protože jsou uloženy na nejvyšší úrovni tabulky, a určovat rámec proměnné. Pokud proměnná není v nejvyšší úrovni, musí být uložena v lokálním rámci, jinak v globálním. Tabulka symbolů také obsahuje jedinečná jména pro každou proměnnou a funkci.